\documentclass[11pt,a4paper]{article}

% ── Packages ─────────────────────────────────────────────────────────────────
\usepackage[utf8]{inputenc}
\usepackage[T1]{fontenc}
\usepackage{lmodern}
\usepackage{microtype}
\usepackage{geometry}
\geometry{margin=2.5cm}

\usepackage{amsmath, amssymb, amsthm}
\usepackage{bm}           % bold math
\usepackage{mathtools}    % dcases, etc.
\usepackage{graphicx}
\usepackage{booktabs}
\usepackage{array}
\usepackage{xcolor}
\usepackage{hyperref}
\hypersetup{
    colorlinks=true,
    linkcolor=blue!60!black,
    citecolor=green!50!black,
    urlcolor=blue!70!black,
}
\usepackage{cleveref}
\usepackage{algorithm}
\usepackage{algpseudocode}
\usepackage{tcolorbox}
\tcbuselibrary{theorems, breakable, skins}

% ── Theorem environments ──────────────────────────────────────────────────────
\newtheorem{theorem}{Theorem}[section]
\newtheorem{proposition}[theorem]{Proposition}
\newtheorem{lemma}[theorem]{Lemma}
\newtheorem{corollary}[theorem]{Corollary}
\newtheorem{definition}[theorem]{Definition}
\newtheorem{remark}[theorem]{Remark}

% ── Custom commands ───────────────────────────────────────────────────────────
\newcommand{\R}{\mathbb{R}}
\newcommand{\N}{\mathbb{N}}
\newcommand{\E}{\mathbb{E}}
\newcommand{\Prob}{\mathbb{P}}
\newcommand{\cS}{\mathcal{S}}
\newcommand{\cA}{\mathcal{A}}
\newcommand{\cO}{\mathcal{O}}
\newcommand{\cT}{\mathcal{T}}
\newcommand{\cR}{\mathcal{R}}
\newcommand{\cZ}{\mathcal{Z}}
\newcommand{\cF}{\mathcal{F}}
\newcommand{\cL}{\mathcal{L}}
\newcommand{\cN}{\mathcal{N}}
\newcommand{\policy}{\pi}
\newcommand{\Vpi}{V^{\policy}}
\newcommand{\Qpi}{Q^{\policy}}
\newcommand{\clip}{\operatorname{clip}}
\newcommand{\diag}{\operatorname{diag}}
\newcommand{\softmax}{\operatorname{softmax}}
\DeclareMathOperator*{\argmax}{argmax}
\DeclareMathOperator*{\argmin}{argmin}

% Coloured box for key results
\newtcolorbox{keyresult}[1][Key Result]{
  colback=blue!5!white, colframe=blue!55!black,
  fonttitle=\bfseries, title=#1, breakable
}

% ─────────────────────────────────────────────────────────────────────────────
\title{%
  \textbf{ANTS: Emergent Ant Colony Simulation}\\[4pt]
  \large Mathematical Foundations and Technical Report\\[2pt]
  \normalsize Phase 1 --- Single-Agent Navigation via Deep Reinforcement Learning
}
\author{%
  Priya\\
  \small OUSUMS Lab, Department of Physics, IISc Bangalore\\
  \small \texttt{priyaa@iisc.ac.in}
}
\date{2025}

% ─────────────────────────────────────────────────────────────────────────────
\begin{document}
\maketitle
\tableofcontents
\vspace{1em}

\begin{abstract}
ANTS (Emergent Ant Colony Simulation) is a deep reinforcement learning project that models
foraging behaviour in ant colonies.
A single agent equipped with only a partial, local view of a procedurally generated maze
must discover a food source and return to the nest.
This document provides the complete mathematical treatment of the problem:
the partially observable Markov decision process (POMDP) formulation,
the potential-based reward shaping theorem and its application to distance-guided learning,
the Proximal Policy Optimisation (PPO) algorithm and its convergence guarantees,
the convolutional neural network architecture used as a feature extractor,
and the frame-stacking technique that converts a POMDP into an approximately Markovian problem.
\end{abstract}

% ═════════════════════════════════════════════════════════════════════════════
\section{Problem Formulation}
% ═════════════════════════════════════════════════════════════════════════════

\subsection{The Foraging Task}

The ant colony foraging problem consists of a single agent navigating an
$W \times H$ grid maze to locate food and return to its nest.
Three features make this a hard learning problem:

\begin{enumerate}
  \item \textbf{Partial observability.} The agent sees only a $(2w{+}1)\times(2w{+}1)$
        window centred on its position (window half-width $w=3$ for a $7\times7$ view).
        In an $11\times11$ maze, this covers at most
        $\frac{7^2}{11^2} \approx 40.3\%$ of cells,
        and typically far less once walls are excluded.

  \item \textbf{Procedural generation.} Each episode uses a freshly generated maze
        (randomised DFS), so the agent cannot memorise a fixed map.
        Successful generalisation requires learning \emph{structural} navigation cues.

  \item \textbf{Sparse primary reward.} A terminal reward of $+10$ is given only on
        food discovery; all intermediate feedback is engineered via reward shaping.
\end{enumerate}

\subsection{Partially Observable Markov Decision Process}

\begin{definition}[POMDP]
A \emph{Partially Observable Markov Decision Process} is a tuple
$\langle \cS, \cA, \cT, \cR, \cO, \cZ, \gamma \rangle$ where:
\begin{itemize}
  \item $\cS$ is the state space;
  \item $\cA$ is the finite action space;
  \item $\cT : \cS \times \cA \to \Delta(\cS)$ is the transition kernel,
        $\cT(s' \mid s, a) = \Prob[S_{t+1} = s' \mid S_t = s, A_t = a]$;
  \item $\cR : \cS \times \cA \times \cS \to \R$ is the reward function;
  \item $\cO$ is the observation space;
  \item $\cZ : \cS \times \cA \to \Delta(\cO)$ is the emission kernel,
        $\cZ(o \mid s, a) = \Prob[O_t = o \mid S_t = s, A_{t-1} = a]$;
  \item $\gamma \in [0,1)$ is the discount factor.
\end{itemize}
\end{definition}

\subsubsection{State Space $\cS$}

The full Markov state at time $t$ is
\begin{equation}
  s_t = \bigl(\mathbf{G},\; \mathbf{p}_t,\; \mathbf{p}_{\text{nest}},\;
              \mathbf{p}_{\text{food}},\; f_t\bigr),
\end{equation}
where:
\begin{itemize}
  \item $\mathbf{G} \in \{0,1\}^{W \times H}$ is the maze grid ($1 = $ wall);
  \item $\mathbf{p}_t = (x_t, y_t) \in \{0,\ldots,W{-}1\} \times \{0,\ldots,H{-}1\}$ is the agent position;
  \item $\mathbf{p}_{\text{nest}}, \mathbf{p}_{\text{food}}$ are fixed entity positions;
  \item $f_t \in \{0,1\}$ is the food-carrying flag.
\end{itemize}

\subsubsection{Action Space $\cA$}

The agent has four movement actions:
\begin{equation}
  \cA = \{{\uparrow}, {\rightarrow}, {\downarrow}, {\leftarrow}\},
  \quad |\cA| = 4.
\end{equation}
Action $a$ maps to displacement $(\Delta x, \Delta y)$:
\begin{equation}
  (\Delta x, \Delta y) =
  \begin{cases}
    (0, -1) & a = 0\ (\uparrow) \\
    (+1, 0) & a = 1\ (\rightarrow) \\
    (0, +1) & a = 2\ (\downarrow) \\
    (-1, 0) & a = 3\ (\leftarrow)
  \end{cases}
\end{equation}

\subsubsection{Transition Kernel $\cT$}

Movement is deterministic:
\begin{equation}
  (x_{t+1}, y_{t+1}) =
  \begin{cases}
    (x_t + \Delta x,\; y_t + \Delta y)
      & \text{if }(x_t{+}\Delta x, y_t{+}\Delta y)\text{ is a valid free cell,} \\
    (x_t, y_t)
      & \text{otherwise (wall or boundary collision).}
  \end{cases}
\end{equation}
A \emph{valid free cell} satisfies $0 \le x < W$, $0 \le y < H$,
and $\mathbf{G}[x,y] = 0$.

\subsubsection{Observation Space $\cO$ and Emission Kernel $\cZ$}

The observation is a three-channel tensor:
\begin{equation}
  o_t \in [0,1]^{3 \times (2w{+}1) \times (2w{+}1)},
  \quad w = 3.
\end{equation}
Channel semantics at local pixel $(i,j)$ (offset $(\Delta x, \Delta y) = (i-w, j-w)$
from the agent):
\begin{align}
  o_t^{(0)}[i,j] &=
  \begin{cases} 1 & \text{if }(x_t{+}\Delta x, y_t{+}\Delta y)\text{ is a wall or OOB,} \\ 0 & \text{otherwise,} \end{cases} \\[4pt]
  o_t^{(1)}[i,j] &=
  \begin{cases} 1 & \text{if }(x_t{+}\Delta x, y_t{+}\Delta y) = \mathbf{p}_{\text{food}} \land \lnot f_t, \\ 0 & \text{otherwise,} \end{cases} \\[4pt]
  o_t^{(2)}[i,j] &=
  \begin{cases} 1 & \text{if }(x_t{+}\Delta x, y_t{+}\Delta y) = \mathbf{p}_{\text{nest}}, \\ 0 & \text{otherwise.} \end{cases}
\end{align}
The emission is deterministic given $s_t$, so $\cZ(o \mid s, a) \in \{0,1\}$.

% ═════════════════════════════════════════════════════════════════════════════
\section{Reward Function}
% ═════════════════════════════════════════════════════════════════════════════

\subsection{Raw Reward Components}

The total reward at each step is the sum of five components:

\begin{equation}
  r_t = r_t^{\text{step}} + r_t^{\text{wall}} + r_t^{\text{dist}} + r_t^{\text{revisit}} + r_t^{\text{food}}.
  \label{eq:reward}
\end{equation}

\begin{table}[h]
\centering
\caption{Reward components and their design rationale.}
\begin{tabular}{@{}llll@{}}
\toprule
Component & Formula & Magnitude & Purpose \\
\midrule
Step penalty & $r^{\text{step}} = -0.01$ & const. & Efficiency pressure \\
Wall penalty & $r^{\text{wall}} = -0.10 \cdot \mathbf{1}[\text{hit wall}]$ & $\le 0$ & Discourage bumping \\
Distance shaping & $r^{\text{dist}} = 0.15 \cdot (d_{t-1} - d_t)$ & $\in[-0.15, +0.15]$ & Guiding gradient \\
Revisitation & $r^{\text{revisit}} = -0.05 \cdot \min(v_t - 1,\; 5)$ & $\in[-0.25, 0]$ & Break oscillation \\
Food reward & $r^{\text{food}} = +10.0 \cdot \mathbf{1}[\text{food reached}]$ & terminal & Primary objective \\
\bottomrule
\end{tabular}
\label{tab:rewards}
\end{table}

Here $d_t = \|\mathbf{p}_t - \mathbf{p}_{\text{food}}\|_1$ is the Manhattan distance
to the food at time $t$, and $v_t \ge 1$ is the visit count for the cell
$\mathbf{p}_t$ in the current episode.

\subsection{The Revisitation Penalty Cap: A Critical Design Choice}

\begin{remark}[Catastrophic penalty explosion]
An early version used $r^{\text{revisit}} = -0.03(v_t - 1)$, \emph{uncapped}.
A cell visited 100 times accumulated $-2.97$ per step,
overwhelming all other signals and producing total episode rewards of $-311.81$.
The agent learned to freeze rather than explore.

The fix caps extra visits at $5$:
\begin{equation}
  r_t^{\text{revisit}} = -0.05 \cdot \min\bigl(v_t - 1,\; 5\bigr),
  \qquad \text{maximum } {-0.25 \text{ per step.}}
\end{equation}
With a step penalty of $-0.01$ and food reward of $+10$, the
maximum undiscounted loss from revisitation over 300 steps is $-75$,
still dominated by the $+10$ terminal reward.
\end{remark}

% ═════════════════════════════════════════════════════════════════════════════
\section{Potential-Based Reward Shaping}
% ═════════════════════════════════════════════════════════════════════════════

\subsection{Policy Invariance Theorem}

The distance shaping term $r^{\text{dist}}$ is an instance of
\emph{potential-based reward shaping} (PBRS), which carries a fundamental optimality guarantee.

\begin{theorem}[Ng, Harada \& Russell, 1999]
\label{thm:pbrs}
Let $M = \langle \cS, \cA, \cT, R, \gamma \rangle$ be an MDP with optimal policy $\policy^*$.
Define the shaped reward
\begin{equation}
  \tilde{R}(s, a, s') = R(s, a, s') + F(s, a, s'),
\end{equation}
where the shaping function $F$ satisfies
\begin{equation}
  F(s, a, s') = \gamma \Phi(s') - \Phi(s)
  \label{eq:pbrs}
\end{equation}
for some \emph{potential function} $\Phi : \cS \to \R$.
Then every optimal policy $\tilde{\policy}^*$ of the shaped MDP
$\tilde{M} = \langle \cS, \cA, \cT, \tilde{R}, \gamma \rangle$
is also optimal in the original MDP $M$.
\end{theorem}

\begin{proof}[Proof sketch]
The shaped value function satisfies
\begin{align}
  \tilde{V}^{\policy}(s)
    &= V^{\policy}(s) - \Phi(s) + \frac{\E[\Phi(s_\infty) \mid s_0=s,\policy]}{(1-\gamma)/\gamma},
\end{align}
where the last term vanishes for episodic tasks with absorbing terminals
(since $\Phi(s_{\text{terminal}}) = 0$ can always be arranged).
Hence $\tilde{V}^{\policy}(s) = V^{\policy}(s) - \Phi(s)$,
and the ordering of policies is preserved:
$\tilde{V}^{\policy_1}(s) \ge \tilde{V}^{\policy_2}(s) \iff V^{\policy_1}(s) \ge V^{\policy_2}(s)$.
\end{proof}

\subsection{Application: Manhattan Distance Potential}

We choose $\Phi(s) = -d(s)$ where $d(s) = \|\mathbf{p} - \mathbf{p}_{\text{food}}\|_1$
is the Manhattan distance from state $s$ to the food.
Then the PBRS formula (\ref{eq:pbrs}) yields
\begin{align}
  F(s_t, a_t, s_{t+1})
    &= \gamma \Phi(s_{t+1}) - \Phi(s_t) \notag \\
    &= -\gamma d_{t+1} + d_t.
\end{align}
Setting $\gamma \approx 1$ and using the approximation $F \approx d_t - d_{t+1}$
gives the implemented form:
\begin{equation}
  r_t^{\text{dist}} = 0.15 \cdot (d_t - d_{t+1}).
  \label{eq:dist_shaping}
\end{equation}
By \cref{thm:pbrs}, this does \emph{not} change the set of optimal policies:
the agent is rewarded for moving toward food but the globally optimal trajectory
(shortest path) remains optimal.

\begin{keyresult}[Shaping bound]
Because $d_t \in [0, W{+}H{-}2]$ and each step changes $d$ by at most $1$,
the shaping signal is bounded:
\[
  r_t^{\text{dist}} \in [-0.15,\; +0.15].
\]
It provides a \emph{dense} gradient toward the food without
overwhelming the sparse $+10$ terminal reward.
\end{keyresult}

% ═════════════════════════════════════════════════════════════════════════════
\section{Proximal Policy Optimisation}
% ═════════════════════════════════════════════════════════════════════════════

\subsection{Policy Gradient Background}

We parameterise the policy as $\policy_\theta(a \mid o)$,
a probability distribution over actions conditioned on the observation,
with parameters $\theta \in \R^d$.
The objective is to maximise the expected discounted return:
\begin{equation}
  J(\theta) = \E_{\tau \sim \policy_\theta}\left[\sum_{t=0}^{T} \gamma^t r_t\right].
\end{equation}
The policy gradient theorem states:
\begin{equation}
  \nabla_\theta J(\theta)
  = \E_{\tau \sim \policy_\theta}\left[
      \sum_{t=0}^{T} \nabla_\theta \log \policy_\theta(a_t \mid o_t) \cdot \hat{A}_t
    \right],
  \label{eq:pg}
\end{equation}
where $\hat{A}_t$ is an estimate of the advantage function
$A^\policy(s_t, a_t) = Q^\policy(s_t, a_t) - V^\policy(s_t)$.

\subsection{Generalised Advantage Estimation}

We use Generalised Advantage Estimation (GAE, Schulman et al.\ 2016)
to reduce variance while controlling bias.
Define the TD residual
\begin{equation}
  \delta_t = r_t + \gamma V_\phi(o_{t+1}) - V_\phi(o_t),
\end{equation}
where $V_\phi$ is a learned value network with parameters $\phi$.
The GAE advantage estimate is
\begin{equation}
  \hat{A}_t^{\text{GAE}(\gamma, \lambda)}
  = \sum_{l=0}^{T-t-1} (\gamma\lambda)^l \delta_{t+l},
  \label{eq:gae}
\end{equation}
with $\lambda \in [0,1]$ controlling the bias--variance tradeoff
($\lambda=0$: one-step TD, low variance, high bias;
 $\lambda=1$: Monte Carlo, high variance, zero bias).
We use $\gamma = 0.99$, $\lambda = 0.95$.

\subsection{PPO Clipped Objective}

Vanilla policy gradient (\ref{eq:pg}) suffers from large, destructive updates.
PPO (Schulman et al.\ 2017) constrains the update implicitly via a clipped ratio:
\begin{equation}
  \rho_t(\theta) = \frac{\policy_\theta(a_t \mid o_t)}{\policy_{\theta_{\text{old}}}(a_t \mid o_t)}.
\end{equation}
The PPO-Clip objective is
\begin{equation}
  \cL^{\text{CLIP}}(\theta)
  = \E_t\left[
      \min\!\left(
        \rho_t(\theta)\,\hat{A}_t,\;\;
        \clip\!\left(\rho_t(\theta),\, 1{-}\epsilon,\, 1{+}\epsilon\right)\hat{A}_t
      \right)
    \right],
  \label{eq:ppo}
\end{equation}
where $\epsilon = 0.2$ is the clip range.

\begin{remark}
The clipping has a natural geometric interpretation.
When $\hat{A}_t > 0$ (a good action), the objective increases $\rho_t$
but is capped at $1+\epsilon$ — the update cannot be too large.
When $\hat{A}_t < 0$ (a bad action), $\rho_t$ is bounded below at $1-\epsilon$.
In both cases, the surrogate is a lower bound on the unclipped objective.
\end{remark}

\subsection{Full PPO Loss}

The total loss combines the policy objective,
a value function regression term, and an entropy bonus:
\begin{equation}
  \cL(\theta, \phi)
  = -\cL^{\text{CLIP}}(\theta)
  + c_v \cL^{\text{VF}}(\phi)
  - c_e \cH[\policy_\theta(\cdot \mid o_t)],
  \label{eq:total_loss}
\end{equation}
where:
\begin{itemize}
  \item $\cL^{\text{VF}}(\phi) = \E_t\left[\bigl(V_\phi(o_t) - \hat{V}_t\bigr)^2\right]$
        is the value function MSE loss;
  \item $\cH[\policy_\theta] = -\sum_a \policy_\theta(a)\log\policy_\theta(a)$
        is the entropy of the current policy;
  \item $c_v = 0.5$ (value coefficient) and $c_e = 0.05$ (entropy coefficient).
\end{itemize}
The entropy bonus $c_e \cH[\policy_\theta]$ discourages premature convergence
to deterministic policies, maintaining exploration pressure throughout training.

\subsection{Training Loop}

\begin{algorithm}[H]
\caption{PPO with Frame-Stacking (Phase 1)}
\label{alg:ppo}
\begin{algorithmic}[1]
\Require Initial parameters $\theta_0, \phi_0$; $N_{\text{envs}}=8$; $T=512$; epochs $K=10$; batch $B=256$
\For{each iteration $i = 1, 2, \ldots$}
  \State Collect $N_{\text{envs}} \times T$ transitions using $\policy_{\theta_{\text{old}}}$
         across 8 parallel environments
  \State Compute advantages $\hat{A}_t^{\text{GAE}}$ via \eqref{eq:gae}
  \State Compute returns $\hat{V}_t = \hat{A}_t + V_\phi(o_t)$
  \For{epoch $k = 1, \ldots, K$}
    \For{each minibatch $\cB$ of size $B$ sampled from the rollout}
      \State Compute $\cL(\theta, \phi)$ via \eqref{eq:total_loss}
      \State Update $\theta \leftarrow \theta - \alpha \nabla_\theta \cL$,
             $\phi \leftarrow \phi - \alpha \nabla_\phi \cL$
    \EndFor
  \EndFor
  \State $\theta_{\text{old}} \leftarrow \theta$
\EndFor
\end{algorithmic}
\end{algorithm}

The effective batch size per update is $N_{\text{envs}} \times T = 4096$ transitions,
processed over $K=10$ epochs in minibatches of $B=256$.

% ═════════════════════════════════════════════════════════════════════════════
\section{Neural Network Architecture}
% ═════════════════════════════════════════════════════════════════════════════

\subsection{Frame Stacking as a Markov Approximation}

A POMDP is in general not Markovian in observations alone.
However, if the history $h_t = (o_{t-N+1}, \ldots, o_t)$ contains enough
information to identify the hidden state, we recover approximate Markovianity.

\begin{proposition}[Frame-stacking Markov approximation]
\label{prop:framestack}
Define the stacked observation
$\bar{o}_t = (o_{t-N+1}, \ldots, o_t) \in \cO^N$.
If the agent's position and velocity can be uniquely inferred from
$N$ consecutive local views, then the process
$(\bar{o}_t, a_t)_{t \ge N}$ satisfies the Markov property.
\end{proposition}

\begin{proof}[Informal argument]
The hidden state relevant to navigation is the agent's position $\mathbf{p}_t$.
Given two consecutive observations $o_{t-1}$ and $o_t$, the shift
between the centred local views reveals the direction of movement,
so $\mathbf{p}_t$ (modulo global maze position) can be inferred.
With $N=4$ frames, direction, speed (always 1 cell), and recent trajectory
are fully determined from $\bar{o}_t$, approximating the Markov condition.
\end{proof}

In practice, frame stacking concatenates $N=4$ observations along the channel axis:
\begin{equation}
  \bar{o}_t \in \R^{(N \cdot C) \times H_{\text{obs}} \times W_{\text{obs}}}
  = \R^{12 \times 7 \times 7},
  \quad N=4,\; C=3.
\end{equation}

\subsection{CNN Feature Extractor}

The policy and value networks share a CNN feature extractor
$\phi_\theta : \R^{12 \times 7 \times 7} \to \R^{64}$:

\begin{equation}
\resizebox{\textwidth}{!}{$
\underbrace{\R^{12 \times 7 \times 7}}_{\bar{o}_t}
\xrightarrow{\text{Conv}_{3\times3}^{12 \to 16} + \text{ReLU}}
\underbrace{\R^{16 \times 7 \times 7}}{}
\xrightarrow{\text{Conv}_{3\times3}^{16 \to 32} + \text{ReLU}}
\underbrace{\R^{32 \times 7 \times 7}}{}
\xrightarrow{\text{AvgPool}(3,3)}
\underbrace{\R^{32 \times 3 \times 3}}{}
\xrightarrow{\text{Flatten}}
\underbrace{\R^{288}}{}
\xrightarrow{\text{Linear}_{288 \to 64} + \text{ReLU}}
\underbrace{\R^{64}}_{\mathbf{z}_t}
$}
\end{equation}

All convolutions use kernel size $3 \times 3$ with padding $1$
(same-padded), preserving spatial dimensions.
The \texttt{AdaptiveAvgPool2d(3, 3)} compresses to $3\times 3$,
yielding $32 \times 9 = 288$ features before the linear head.

\subsubsection{Why Same-Padding?}

With input size $7 \times 7$, padding $= 1$ and kernel $= 3$, the output is:
\begin{equation}
  H_{\text{out}} = \left\lfloor \frac{H_{\text{in}} + 2p - k}{s} + 1 \right\rfloor
  = \left\lfloor \frac{7 + 2 - 3}{1} + 1 \right\rfloor = 7.
\end{equation}
Same-padding preserves spatial structure so that each position in the feature map
still corresponds to a specific region in the local observation window.
This matters: the model must localise food and walls \emph{relative to the agent}.

\subsubsection{Parameter Count}

\begin{table}[h]
\centering
\caption{CNN parameter count.}
\begin{tabular}{@{}lrrr@{}}
\toprule
Layer & Output shape & Parameters & Notes \\
\midrule
Conv1 (k=3, p=1) & $16\times7\times7$ & $12\times3\times3\times16 + 16 = 1{,}744$ & \\
Conv2 (k=3, p=1) & $32\times7\times7$ & $16\times3\times3\times32 + 32 = 4{,}640$ & \\
AvgPool(3,3)     & $32\times3\times3$ & $0$ & \\
Flatten          & $288$              & $0$ & \\
Linear + ReLU    & $64$               & $288\times64 + 64 = 18{,}496$ & \\
\midrule
\textbf{Total CNN} & & \textbf{24,880} & \\
\bottomrule
\end{tabular}
\end{table}

The small parameter count (under 25k) is intentional:
the $7\times7$ local view contains limited information,
and a larger extractor would overfit the maze topology rather than
learning generalizable navigation heuristics.

\subsubsection{Policy and Value Heads}

The feature vector $\mathbf{z}_t \in \R^{64}$ feeds two heads:
\begin{align}
  \policy_\theta(a \mid \bar{o}_t) &= \softmax\!\left(\mathbf{W}_\pi \mathbf{z}_t + \mathbf{b}_\pi\right),
    \quad \mathbf{W}_\pi \in \R^{4 \times 64}, \\
  V_\phi(\bar{o}_t) &= \mathbf{w}_v^\top \mathbf{z}_t + b_v,
    \quad \mathbf{w}_v \in \R^{64}.
\end{align}

% ═════════════════════════════════════════════════════════════════════════════
\section{Maze Generation}
% ═════════════════════════════════════════════════════════════════════════════

\subsection{Randomised Depth-First Search}

Each episode generates a new \emph{perfect maze} (exactly one path between
any two cells) via the randomised DFS algorithm:

\begin{algorithm}[H]
\caption{Randomised DFS Maze Generation}
\label{alg:dfs}
\begin{algorithmic}[1]
\State Initialise all cells as walls: $\mathbf{G} \leftarrow \mathbf{1}^{W \times H}$
\State Push starting cell $c_0$ onto stack $\mathcal{K}$; mark $c_0$ as visited
\While{$\mathcal{K}$ is non-empty}
  \State $c \leftarrow \mathcal{K}.\text{top}()$
  \State $N \leftarrow$ unvisited neighbours of $c$ at distance 2 (grid step)
  \If{$N \ne \emptyset$}
    \State Choose $n \in N$ uniformly at random
    \State Carve passage: set $\mathbf{G}[c] = \mathbf{G}[\text{midpoint}(c,n)] = \mathbf{G}[n] = 0$
    \State Mark $n$ as visited; push $n$ onto $\mathcal{K}$
  \Else
    \State $\mathcal{K}.\text{pop}()$
  \EndIf
\EndWhile
\end{algorithmic}
\end{algorithm}

The resulting maze has exactly $(WH - 1)/4$ passages (for odd $W, H$),
guaranteeing connectedness and the perfect-maze property.

\subsection{Entity Placement}

Nest and food positions are sampled uniformly from the set of free cells
subject to a minimum Manhattan distance:
\begin{equation}
  \|\mathbf{p}_{\text{nest}} - \mathbf{p}_{\text{food}}\|_1 \ge d_{\min} = 3.
\end{equation}
This prevents degenerate episodes where the food is immediately visible
in the starting observation window.

% ═════════════════════════════════════════════════════════════════════════════
\section{Phase 2 Preview: Multi-Agent Extensions}
% ═════════════════════════════════════════════════════════════════════════════

\subsection{Pheromone Diffusion}

Phase 2 introduces stigmergic communication via a pheromone field
$\mathbf{P} \in \R^{n_c \times W \times H}_{\ge 0}$ with $n_c$ channels.
The physics are governed by two processes:

\paragraph{Evaporation.}
\begin{equation}
  \mathbf{P}' \leftarrow \mathbf{P} \cdot (1 - \varepsilon),
  \quad \varepsilon \in (0,1).
\end{equation}

\paragraph{Diffusion.}
\begin{equation}
  \mathbf{P}^{(k)} \leftarrow \mathbf{P}^{\prime(k)} * \mathbf{K},
  \quad k = 1, \ldots, n_c,
\end{equation}
where $*$ denotes 2-D convolution and $\mathbf{K}$ is a normalised Gaussian kernel
(e.g.\ $3\times3$ with $\sigma = 0.7$, $\sum_{ij} K_{ij} = 1$).

Combined update per step:
\begin{equation}
  \mathbf{P}_{t+1} = \clip\!\left(
    \left(\mathbf{P}_t \cdot (1-\varepsilon)\right) * \mathbf{K},
    \; 0, \; P_{\max}
  \right).
  \label{eq:pheromone}
\end{equation}

\subsection{Centralised Training, Decentralised Execution (CTDE)}

With $n$ agents, each agent $i$ observes $o_t^{(i)}$ locally.
At training time a centralised critic receives the joint state
$\mathbf{s}_t = (s_t^{(1)}, \ldots, s_t^{(n)})$
and pheromone field $\mathbf{P}_t$:
\begin{equation}
  V^\pi(\mathbf{s}_t, \mathbf{P}_t) \approx V_\phi\!\left(
    f_\psi(o_t^{(1)}, \ldots, o_t^{(n)}, \mathbf{P}_t)
  \right),
\end{equation}
while at execution time each agent acts with
$\policy^{(i)}_\theta(a \mid o_t^{(i)}, \mathbf{P}_t^{\text{local}})$.
This is the MAPPO framework (Yu et al.\ 2022).

% ═════════════════════════════════════════════════════════════════════════════
\section{Empirical Results (Phase 1)}
% ═════════════════════════════════════════════════════════════════════════════

\begin{table}[h]
\centering
\caption{Phase 1 training configuration and key results.}
\begin{tabular}{@{}ll@{}}
\toprule
\textbf{Setting} & \textbf{Value} \\
\midrule
Algorithm & PPO + Frame Stacking ($N=4$) \\
Maze size & $11 \times 11$ (procedurally generated) \\
Observation & $7 \times 7$ local window ($3$ channels) \\
Stacked obs. shape & $(12, 7, 7)$ \\
Parallel envs. & $8$ \\
Total training steps & $500{,}000$ \\
Learning rate & $3 \times 10^{-4}$ \\
Rollout length $T$ & $512$ \\
Mini-batch size $B$ & $256$ \\
PPO epochs $K$ & $10$ \\
Entropy coefficient $c_e$ & $0.05$ \\
Clip range $\epsilon$ & $0.2$ \\
Discount $\gamma$ & $0.99$ \\
GAE $\lambda$ & $0.95$ \\
\midrule
\textbf{Best episode} & \textbf{16 steps to food} \\
Reward at convergence & $\approx 9.64$ \\
Observation coverage & $\le 40.3\%$ of maze \\
\bottomrule
\end{tabular}
\end{table}

\begin{keyresult}[Spatial reasoning under partial observability]
The trained agent solves randomly generated $11\times11$ mazes in as few as
\textbf{16 steps} using only a $7\times7$ local observation window —
covering less than 40\% of the maze.
The agent demonstrates genuine spatial reasoning:
identifying corridor junctions, navigating toward the food signal,
and avoiding oscillation loops via the revisitation penalty.
\end{keyresult}

% ═════════════════════════════════════════════════════════════════════════════
\section{Discussion and Future Work}
% ═════════════════════════════════════════════════════════════════════════════

\subsection{Why Frame Stacking Outperforms LSTM Here}

RecurrentPPO with LSTM requires gradients to flow back through time,
making training sensitive to:
(i) rollout length (too short $\Rightarrow$ gradients vanish before reaching
    informative transitions),
(ii) reward noise (noisy gradients in the BPTT window destabilise
    the cell state), and
(iii) wall-clock time ($\approx 3\times$ slower per step due to sequential
    hidden state propagation).

Frame stacking avoids all three issues by treating temporal context
as a spatial problem: $N$ observations are concatenated into extra channels,
so the CNN processes them in parallel.
The tradeoff is that frame stacking provides only a fixed-length history,
while LSTM maintains a theoretically infinite context.
For the maze navigation task, $N=4$ frames are sufficient
(velocity is determined from 2 frames; $N=4$ also captures
recent turning decisions), so LSTM's added complexity is not justified.

\subsection{Curriculum Learning Roadmap}

The \texttt{EpisodicMemory} module implements automatic difficulty scaling:
\begin{equation}
  d_{\min}^{(t+1)} =
  \begin{cases}
    d_{\min}^{(t)} + 1 & \text{if success rate} > 0.75 \\
    d_{\min}^{(t)} - 1 & \text{if success rate} < 0.25 \\
    d_{\min}^{(t)}     & \text{otherwise.}
  \end{cases}
\end{equation}
This adaptive curriculum ensures the task difficulty tracks the agent's
current competence level.

\subsection{Emergent Communication (Phase 3)}

Phase 3 will investigate whether agents develop interpretable
communication protocols by training with a discrete message channel
$m_t^{(i)} \in \{0,\ldots,M-1\}^d$ alongside pheromone signals.
Interpretability will be assessed via linear probing
(Alain \& Bengio, 2017): a logistic regression trained on agent messages
as features to predict latent state variables (food direction, carry flag, etc.).
Non-trivial probing accuracy would indicate that the learned protocol
encodes semantically meaningful information.

% ═════════════════════════════════════════════════════════════════════════════
\section*{References}
% ═════════════════════════════════════════════════════════════════════════════

\begin{enumerate}
  \item Ng, A.~Y., Harada, D., \& Russell, S.~J.\ (1999).
        \textit{Policy invariance under reward transformations: Theory and application
        to reward shaping.}
        Proceedings of the 16th ICML.

  \item Schulman, J., Moritz, P., Levine, S., Jordan, M., \& Abbeel, P.\ (2016).
        \textit{High-dimensional continuous control using generalized advantage estimation.}
        ICLR 2016.

  \item Schulman, J., Wolski, F., Dhariwal, P., Radford, A., \& Klimov, O.\ (2017).
        \textit{Proximal Policy Optimization Algorithms.}
        arXiv:1707.06347.

  \item Mnih, V., et al.\ (2015).
        \textit{Human-level control through deep reinforcement learning.}
        \textit{Nature}, 518(7540), 529--533.

  \item Yu, C., Velu, A., Vinitsky, E., Gao, J., Wang, Y., Bayen, A., \& Wu, Y.\ (2022).
        \textit{The surprising effectiveness of PPO in cooperative, multi-agent games.}
        NeurIPS 2022.

  \item Alain, G., \& Bengio, Y.\ (2017).
        \textit{Understanding intermediate layers using linear classifier probes.}
        ICLR 2017 Workshop.

  \item Bonabeau, E., Dorigo, M., \& Theraulaz, G.\ (1999).
        \textit{Swarm Intelligence: From Natural to Artificial Systems.}
        Oxford University Press.
\end{enumerate}

\end{document}
